\begin{table}[htbp]
\centering
\caption{Every reported quantity with and without the HN2 vial, whose contents are not established (Section~S5).}
\label{tab:si-hn2}
\small
\begin{tabular}{lrr}
\toprule
Quantity & All acquisitions & HN2 excluded \\
\midrule
Queries & 30 & 27 \\
Measured, molecular & 30 & 27 \\
Computed, molecular (11-member) & 13 & 13 \\
Computed, molecular (54-member) & 13 & 13 \\
Measured, class & 30 & 27 \\
Computed, class (11-member) & 25 & 25 \\
Computed, class (54-member) & 25 & 25 \\
Measured self-similarity & 0.983 & 0.981 \\
Computed self-similarity & 0.674 & 0.742 \\
Measured margin & 0.219 & 0.150 \\
Computed margin & 0.056 & 0.052 \\
\bottomrule
\end{tabular}

\begin{flushleft}\footnotesize Excluding HN2 raises the computed-library self-similarity from 0.674 to 0.742, so the deficit against the measured library falls from 0.309 to 0.239, and it removes the entire class-level difference between the two libraries. The HN2 acquisitions also carry an atypically large measured margin, which inflates the measured mean. The main text quotes the full-set values, with the HN2-excluded values given alongside at each headline; Section~\ref{si:hn2sens} gives the comparison in full.\end{flushleft}
\end{table}
